% Generated draft: requires mathematical and source review.
\documentclass[11pt]{article}
\usepackage[margin=25mm]{geometry}
\usepackage{fontspec}
\setmainfont{Latin Modern Roman}
\usepackage{amsmath,amssymb,mathtools,unicode-math}
\setmathfont{Latin Modern Math}
\usepackage{enumitem,needspace,xurl,hyperref}
\setlength{\parindent}{0pt}
\setlength{\parskip}{5pt}
\setlength{\emergencystretch}{3em}
\newcommand{\N}{\mathbb{N}}
\newcommand{\Z}{\mathbb{Z}}
\newcommand{\Q}{\mathbb{Q}}
\newcommand{\R}{\mathbb{R}}
\newcommand{\C}{\mathbb{C}}
\newcommand{\F}{\mathbb{F}}
\newcommand{\A}{\mathbb{A}}
\newcommand{\PP}{\mathbb{P}}
\newcommand{\E}{\mathbb{E}}
\newcommand{\eps}{\varepsilon}
\newcommand{\dd}{\,\mathrm{d}}
\begin{document}
\section*{UnsolvedMath ID: 30000957}
\textbf{Source problem code:} catalogue entry
\subsection{Definitions and mathematical statement}
In the usual algebraic interpretation, for a connected manifold $N$ and complex semisimple algebraic group $G$, its character variety is
\[
X_G(N)=\operatorname{Hom}(\pi_1(N),G)//G,
\]
where the quotient is by conjugation. Boundary inclusion induces restriction $\rho_M:X_G(M)\to X_G(F)$ for $F=\partial M$; for disconnected boundaries use the corresponding componentwise character varieties. The requested target is $\operatorname{im}(\rho_M)$.

For the compact oriented three-manifold and Witten--Reshetikhin--Turaev theory conventions intended by the source, let $V_{G,r}(F)$ be the boundary state space at admissible level $r$ and $I_{G,r}(M)$ the vector assigned to the filling. The conjecture, identified as Sikora's Conjecture 6, asks whether the entire collection
\[
\{I_{G,r}(M)\in V_{G,r}(F):r\text{ an admissible level}\}
\]
determines the boundary character-variety image. In particular, for two fillings $M,N$ with the same identified boundary and consistent quantum choices, it requires
\[
\left[\forall r,\ I_{G,r}(M)=I_{G,r}(N)\right]
\quad\Longrightarrow\quad
\operatorname{im}(\rho_M)=\operatorname{im}(\rho_N).
\]
This expresses determination as uniqueness from all levelwise vectors; the source does not give a reconstruction algorithm.

The extract does not specify compactness/orientation, the algebraic versus compact-real-form convention for $G$, allowed levels and quantum normalization, or boundary parametrizations and anomaly/framing data. It also does not say whether ``image'' means the actual constructible image, its Zariski closure or a scheme-theoretic image. These targets are distinguished rather than silently replacing the literal image by its closure. The conjecture concerns vectors in fixed boundary spaces over all levels, not only scalar closed-manifold invariants or finitely many levels.
\subsection{Short English statement}
Do all WRT boundary vectors of a three-manifold determine which boundary characters extend to it? The exact image notion and the geometric and quantum conventions remain unspecified; no finite-level or scalar-invariant replacement is made.
\subsection{Sources}
\begin{itemize}
\item Supplied report: \emph{Invariants in Low-Dimensional Topology} (2009), Oberwolfach Reports, DOI 10.4171/owr/2008/22, pp. 1157--1232; proposed year 2008. \url{https://doi.org/10.4171/owr/2008/22}.
\item Supplied Sikora 2008 reference: \url{https://arxiv.org/abs/0807.0943}; the record identifies Conjecture 6.
\item Local UnsolvedMath registry, permanent ID 30000957, catalogue entry; created 25 August 2026, updated 16 September 2026. Its corrected\_verified note calls this the complete canonical extraction of a repeated passage. The supplied 21 August 2026 triage labels the determination conjecture open, distinguishing construction of a quantum character-variety algebra from proof that all WRT vectors determine the boundary image. This dated assessment has not been independently verified here.
\end{itemize}
\end{document}
